\section{Common factors of the highest-degree kernel minors}
\label{sec:digit-kernel-divisors}

A scalar factor common to every tensor coordinate may be divided out
before a local reconstruction is attempted. The following calculation
identifies such factors in the highest power of the normalizer in
Section~\ref{sec:deformed-crossing}. It explains why the order at the
point $X=1$ can be much larger than the orders at the other roots of
$1+X^{2(q-1)}$. Dividing these factors does not, by itself, bound the
ranks of the remaining coefficients.

Work over $\F_2$. Fix $q=2^p$, $p\geq1$, and $s\geq1$, and put
\[
 M=2^s,\qquad N=q-1,\qquad
 H=\frac{q^s-1}{q-1},\qquad H'=\frac{q^{s-1}-1}{q-1},
 \qquad y=X^{2N}.
\]
For $0\leq i<M$ let $E(i)=\sum_{j=0}^{s-1}\operatorname{bit}_j(i)q^j$.
Besides the bitwise intersection $i\mathbin{\&}j$, use
$i\mathbin{|}j$ for bitwise union. Define $M$-by-$M$ matrices
\[
 C_{ij}=y^{E(i\mathbin{\&}j)},\qquad
 U_{ij}=y^{E(i\mathbin{|}j)},\qquad
 \mathcal K=\begin{pmatrix}C&XU\\XU&C\end{pmatrix}.
\]
The greatest common divisor below is monic and is taken in $\F_2[X]$.

\begin{proposition}[The middle minors of the digit kernel]
\label{prop:digit-kernel-minor-divisor}
The greatest common divisor of all $M$-by-$M$ minors of $\mathcal K$ is
\begin{equation}\label{eq:digit-kernel-gcd}
 G=(1+y)^{MH'/2}(1+X)^{M q^{s-1}}.
\end{equation}
In the deformed crossing determinant with $n=M/2$ indices in every
row and column group, every coordinate of the coefficient of $f^{2n}$
is divisible by $G^2$. This remains true after restricting the symbols
to any prescribed equal-energy class. No equality is asserted for the
greatest common divisor after that restriction.
\end{proposition}

\begin{proof}
In binary tensor coordinates,
\[
 C=\bigotimes_{j=0}^{s-1}
       \begin{pmatrix}1&1\\1&y^{q^j}\end{pmatrix},\qquad
 J=\bigotimes_{j=0}^{s-1}
       \begin{pmatrix}0&y^{q^j}\\1&0\end{pmatrix}.
\]
A direct two-by-two multiplication gives $U=CJ$,
$J^{\mathsf T}C=CJ$, and $J^2=y^H I$.
The following row and column operations are polynomially unimodular:
\begin{equation}\label{eq:digit-kernel-congruence}
 \begin{pmatrix}I&0\\XJ^{\mathsf T}&I\end{pmatrix}
 \mathcal K
 \begin{pmatrix}I&XJ\\0&I\end{pmatrix}
 =\begin{pmatrix}C&0\\0&(1+X^2y^H)C\end{pmatrix}.
\end{equation}
Indeed $J^{\mathsf T}C=CJ$ cancels the lower-left block, and
$J^2=y^H I$ gives the lower-right block. Since $1+NH=q^s$,
\[
 1+X^2y^H=(1+X)^{2q^s}.
\]

We compute valuations of determinantal divisors at each irreducible
polynomial. Invertible row and column operations over the corresponding
local ring preserve the ideal generated by minors of a fixed size:
Cauchy--Binet expresses each transformed minor as a linear combination
of the old ones, and the inverse operations give the reverse inclusion.
For a diagonal matrix, the smallest valuation of an $r$-minor is the
sum of its $r$ smallest diagonal valuations. These two elementary facts
suffice for the calculation.

Adding the first row to the second and then the first column to the
second changes each factor of $C$ into
$\operatorname{diag}(1,1+y^{q^j})$.
Since $q$ is a power of two, $1+y^{q^j}=(1+y)^{q^j}$.
Thus $C$ is polynomially equivalent to the diagonal matrix with entries
$(1+y)^{E(i)}$ for $0\leq i<M$.

Every irreducible factor of $1+X^N$ occurs with multiplicity one,
since $N$ is odd. Hence every factor of $1+y=(1+X^N)^2$ has
multiplicity two. At a factor other than $X+1$, the diagonal valuations
of $\mathcal K$ are the numbers $2E(i)$, each repeated twice.
The function $E(i)$ is strictly increasing in $i$: at the highest
changed bit, its contribution exceeds the sum of all lower bit weights.
The $M$ smallest valuations therefore sum to
\[
 4\sum_{i=0}^{M/2-1}E(i)=MH'.
\]
Here each of the $s-1$ lower bits is present in exactly half of those
$M/2$ integers; the sum is zero when $s=1$.

At $X+1$ the valuations are instead
\[
 \{2E(i):0\leq i<M\}
 \ \cup\ \{2E(i)+2q^s:0\leq i<M\}.
\]
Since $H<q^s$, all entries of the first group are smaller than all
entries of the second. Their sum is
$2\sum_{i=0}^{M-1}E(i)=MH$.
There are no other factors: the displayed equivalent matrices have
unit determinants at all other nonzero primes, and $\mathcal K(0)$ is
invertible because $C(0)$ is a tensor product of invertible matrices
$\left(\begin{smallmatrix}1&1\\1&0\end{smallmatrix}\right)$.
The polynomial in \eqref{eq:digit-kernel-gcd} has valuation $MH'$
at every other factor of $1+y$, and valuation
$MH'+M q^{s-1}=MH$ at $X+1$. It is therefore the exact greatest
common divisor.

For the last assertion, recall that the entries involving the $b$ row
or column group carry one factor $f$, except for the $bd$ and $db$
blocks. A term of degree $2n$ in $f$ must use all $n$ such row factors
and all $n$ such column factors. It consequently uses no entry of
those two exceptional blocks. The $a,c$ groups together have $2n$
rows and columns, exactly the number needed to match the $d,b$ groups.
Thus no $ac$ or $ca$ entry can occur either. After grouping $a,c$
against $d,b$, the highest coefficient is the determinant of a
bipartite block matrix, hence the product of two $2n$-minors of
$\mathcal K$. In characteristic two there is no sign to keep track of.
For $2n=M$, each minor is divisible by $G$, proving the claim for
all index choices and therefore for every restricted symbol class.
\end{proof}

For $q=4$, $s=2$, one has $M=4$, $H=5$ and $H'=1$, so
\[
 G^2=(1+X^6)^4(1+X)^{32}.
\]
Its degree is $56$; its order at $X=1$ is $40$, while its order at
each other root of $1+X^3$ is only $8$.
This difference between the distinguished root and the other roots
must be preserved in any degree estimate after division.
The program \emph{Digit-kernel determinantal divisors} in the
\href{\workbenchurl}{workbench} computes every middle minor in the
small alphabets and verifies the exact greatest common divisor.
For larger alphabets the proposition supplies the divisor directly;
an improving construction would still need an adjacent-cut rank
bound for the divided coefficient tensors.

\subsection{The crossing recovered after division}

The common factor can be removed without losing the desired output.
The resulting polynomial deformation still requires a coefficient-rank
estimate, but its value at the distinguished point is exact.
Let $n=M/2$, and let $\mathcal S$ be the set of all $n$-element
subsets of $\{0,\ldots,M-1\}$. For $A,B,C_0,D\in\mathcal S$ put
\[
 H_X(A,B,C_0,D)=
 \frac{\det\mathcal K[A\cup(M+C_0),\,B\cup(M+D)]}{G},
 \qquad
 \tau(A)=\{M-1-i:i\notin A\}.
\]
Here $M+C_0=\{M+i:i\in C_0\}$, and the indices in each block
are placed in increasing order. The notation $C_0$ denotes a subset,
not the matrix $C$. Proposition~\ref{prop:digit-kernel-minor-divisor}
makes every entry of $H_X$ a polynomial.

\begin{proposition}[A polynomial crossing after removing the common factor]
\label{prop:normalized-minor-crossing}
For all $A,B,C_0,D\in\mathcal S$,
\begin{equation}\label{eq:normalized-minor-output}
 H_1(A,B,C_0,D)=[C_0=\tau(A)]\,[D=\tau(B)].
\end{equation}
The map $\tau$ is an involution preserving $\sum_{i\in A}E(i)$.
Thus the same statement gives a crossing after restriction to any
$\tau$-stable family, including every complete fixed-energy class.
The tensor $H_0$ on the full subset spaces has rank one across
$AB\mid C_0D$; on any restricted family its rank is at most one.

Over $\F_2(X)$ the full tensor $H_X$ is locally equivalent to
the balanced middle compound of
\begin{equation}\label{eq:scalar-balanced-kernel}
 \begin{pmatrix}I&XI\\Xy^H I&I\end{pmatrix}.
\end{equation}
These local changes depend on $X$ and do not assert equality of
coefficient ranks.

Finally, let $A_{2n}(X)$ be the coefficient of $f^{2n}$ in the
four-group deformed determinant, with every symbol $(I,J)$ consisting
of two $n$-subsets in the stated finite alphabet. On every complete
fixed combined-energy symbol class, $A_{2n}/G^2$ is a polynomial
crossing deformation and
\begin{equation}\label{eq:normalized-top-output}
 (A_{2n}/G^2)(1;s_a,s_b,s_c,s_d)
   =[s_c=\tau_2(s_a)]\,[s_d=\tau_2(s_b)],
 \qquad \tau_2(I,J)=(\tau(I),\tau(J)).
\end{equation}
Its value at zero has rank at most one across $ad\mid bc$.
\end{proposition}

\begin{proof}
Write $L$ and $R$ for the two unimodular matrices on the left
of \eqref{eq:digit-kernel-congruence}. In characteristic two they
are their own inverses. Consequently
\[
 \mathcal K=L\,\operatorname{diag}(C,hC)\,R,
 \qquad h=(1+X)^{2q^s}.
\]
At $X+1$, the diagonal valuations of $C$ are the $M$ numbers
$2E(i)$, all at most $2H$, and their sum is $MH$.
Any nonzero middle minor of $\operatorname{diag}(C,hC)$ other
than the full first block uses some $j\geq1$ rows and columns
from its second block. Removing $j$ entries from the first block
can decrease the valuation by at most $2jH$, whereas the second
block contributes at least $2jq^s$. Its valuation is therefore
at least
\[
 MH+2j(q^s-H)>MH.
\]
This argument is valid before diagonalization as well: invertible
row and column operations over the local ring preserve all ideals
of minors, as proved above. Since $G$ has valuation $MH$, every
such entry of the middle compound divided by $G$ vanishes at one.
The remaining entry has value one, because
\[
 \frac{\det C}{G}
   =\left(\frac{1+X^N}{1+X}\right)^{M q^{s-1}},
 \qquad
 \left.\frac{1+X^N}{1+X}\right|_{X=1}=N=1\quad\text{in }\F_2.
\]

Apply Cauchy--Binet to the middle compound of the displayed
factorization. After division by $G$ and evaluation at one,
only the first-block intermediate index survives. The answer
is the outer product of the maximal minors of
\[
 \begin{pmatrix}I\\J(1)^{\mathsf T}\end{pmatrix}
 \quad\text{and}\quad
 \begin{pmatrix}I&J(1)\end{pmatrix}.
\]
The matrix $J(1)$ is the permutation $i\mapsto M-1-i$.
The balanced row minor is therefore one precisely when
$C_0=\tau(A)$, and the balanced column minor is one precisely
when $D=\tau(B)$. This proves \eqref{eq:normalized-minor-output}.

Bit complementation gives $E(M-1-i)=H-E(i)$, while
$\sum_{i=0}^{M-1}E(i)=MH/2=nH$. Thus
\[
 \sum_{i\in\tau(A)}E(i)
 =nH-\sum_{i\notin A}E(i)
 =\sum_{i\in A}E(i).
\]
The definition also gives $\tau^2(A)=A$.
At zero, $G(0)=1$ and
$\mathcal K(0)=\operatorname{diag}(C(0),C(0))$.
The balanced minor factors as
$\det C(0)[A,B]\det C(0)[C_0,D]$.
Since $C(0)$ is invertible, its $n$th compound is nonzero,
proving the full rank-one assertion; restriction can only lower it.

For the generic local equivalence, multiply $\mathcal K$ on the
left by $\operatorname{diag}(C^{-1},C^{-1})$ to obtain
$\left(\begin{smallmatrix}I&XJ\\XJ&I\end{smallmatrix}\right)$.
Multiply this matrix on the right by
$\operatorname{diag}(I,J^{-1})$ and on the left by
$\operatorname{diag}(I,J)$. The result is
\eqref{eq:scalar-balanced-kernel}, because $J^2=y^H I$.
All three operations are invertible over $\F_2(X)$ and act
within the separate mode blocks. Their exterior powers are
therefore local invertible maps on the four subset spaces.
The scalar $G^{-1}$ can also be absorbed in one such map.
Nothing in this reasoning commutes coefficient extraction
with an $X$-dependent map.

To check the last assertion with the original party ordering,
the highest coefficient is the product of the two middle minors
with row and column indices
\[
 (I_a,I_c;J_d,J_b),\qquad (I_d,I_b;J_a,J_c).
\]
Equation~\eqref{eq:normalized-minor-output} makes their divided
values at one equal to
\[
 [I_c=\tau(I_a)][J_b=\tau(J_d)]\,
 [I_b=\tau(I_d)][J_c=\tau(J_a)].
\]
Using $\tau^2=1$ gives exactly
\eqref{eq:normalized-top-output}. The preserved subset energies
make $\tau_2$ a permutation of every complete symbol class in
the statement. At zero the first minor splits between $(a,d)$
and $(c,b)$, and the second between $(d,a)$ and $(b,c)$, proving
the stated adjacent-cut rank bound. This remains true after
restricting each local symbol space.
\end{proof}

The proposition identifies the output and a regular polynomial
family. It does not transfer a rank bound from the scalar family
through the parameter-dependent local maps. The unresolved issue
is the rank of the divided digit coefficients, or of suitable
linear combinations of them, across the two adjacent cuts.
